
\section{Simulation: a balanced three-factor economy}
\label{sec:balanced_simulation}
We study direct portfolio learning in a stock-level economy with three economic
factors. The managed spectrum is generated by the characteristic distribution
and the Mat\'ern kernel. Numerical kernel rank is distinct from the number of
economic factors. Table~\ref{tab:sim_calibration} gives the frozen calibration.

\subsection{Characteristics, returns, and the population policy}
There are $N=3G$ stocks and $D=6$ characteristics. Independent group-level states
start from $S_{g,d,0}\sim N(0,1)$ and evolve as
\[
S_{g,d,t+1}=\rho_d S_{g,d,t}+\sqrt{1-\rho_d^2}\,\eta_{g,d,t+1},
\qquad \eta_{g,d,t+1}\sim N(0,1).
\]
With $U_{g,d,t}=\Phi(S_{g,d,t})$, the three roles in group $g$ have
\[
Z_{g,r,d,t}=2\bigl((U_{g,d,t}+r/3)\bmod 1\bigr)-1,
\qquad r=0,1,2.
\]
The same phase shift applies to every coordinate. These population ranks have a
continuous uniform distribution on the cube; they are not empirical ranks.
Writing $u(z)=(z_1+1)/2$, define
\[
\beta(z)=\sqrt{c_\beta}
\begin{pmatrix}
\sqrt{2/3}\cos(2\pi u(z))\\
\sqrt{2/3}\sin(2\pi u(z))\\
1/\sqrt 3
\end{pmatrix}.
\]
The trigonometric identities for three equally spaced phases imply
$\|\beta(z)\|^2=c_\beta$ and
$B_t'B_t/N=\Gamma_\beta=(c_\beta/3)I_3$ at every date.
Independent next-period shocks satisfy
\[
F_{t+1}\sim N(\mu_F,I_3-\mu_F\mu_F'),\qquad
\varepsilon_{t+1}\sim N(0,\sigma_\varepsilon^2 I_N),\qquad
R_{t+1}=B_tF_{t+1}+\varepsilon_{t+1}.
\]
Thus $\mathbb E[F_{t+1}F_{t+1}']=I_3$ is a second-moment normalization.
The conditional mean and second moment of stock returns are
$m_t=B_t\mu_F$ and $S_t=B_tB_t'+\sigma_\varepsilon^2I_N$.

For the paper's weights $w_t(W)=W(Z_t)/N$, minimizing
$\mathbb E[(1-w_t(W)'R_{t+1})^2\mid\mathcal F_t]$
gives $S_tw_t^\star=m_t$. The identity
\[
(BB'+\sigma^2 I)^{-1}B=B(B'B+\sigma^2 I_3)^{-1}
\]
and exact cross-sectional balance therefore yield
\[
w_t^\star=\frac1N B_t
\left(\Gamma_\beta+\frac{\sigma_\varepsilon^2}{N}I_3\right)^{-1}\mu_F,
\qquad
W^\star(z)=\frac{\beta(z)'\mu_F}{c_\beta/3+\sigma_\varepsilon^2/N}.
\]
The population policy is derived from the return model and the loss; it is not
an independently imposed simulation target. With $g=c_\beta/3$,
\[
q^\star=\frac{g}{g+\sigma_\varepsilon^2/N}\|\mu_F\|^2,\quad
Q(W^\star)=1-q^\star,\quad
(\operatorname{SR}^\star)^2=\frac{q^\star}{1-q^\star}.
\]

\subsection{Assumption audit and spectrum}
Stationary initialization and stable Gaussian AR dynamics give geometric
beta-mixing; measurable rank and phase maps and independent next-period shocks
preserve this property (E1). For $A_t=K_t/N^2$, the exact noncentral Gaussian
quadratic-form MGF verifies
\[
\log\mathbb E\!\left[\exp\left(R_{t+1}'A_tR_{t+1}/B_X^2\right)
\mid\mathcal F_t\right]\leq\tfrac12,
\qquad B_X^2=4(c_\beta+\sigma_\varepsilon^2).
\]
Conditional scalar Gaussian moments give $K_4=3$ (E2). The explicit sine,
cosine, and constant loading functions have smooth extensions on the bounded
cube. They belong to the Sobolev space of order $\nu+D/2=9/2$, equivalent to
the Mat\'ern-$3/2$ RKHS, as does their fixed linear combination $W^\star$ (E3).
The conditional pricing restriction follows from
$v_t'(m_t-S_tw_t^\star)=0$ for every admissible $v_t$ (E4).

For the uncentered managed second-moment operator,
\[
\frac{\sigma_\varepsilon^2}{N}T_K\preceq\Sigma_{\mathcal H_K}
\preceq\left(c_\beta+\frac{\sigma_\varepsilon^2}{N}\right)T_K.
\]
The lower term is the idiosyncratic contribution; the upper bound follows by
Jensen's inequality and $\|\beta\|^2=c_\beta$. Thus the managed spectrum has
the Mat\'ern exponent $b=1+2\nu/D=3/2$ (E5). The finite-sample audit uses one
independent population quadrature with 3,072 nodes, including the within-triplet
dependence, and tests the Loewner and ordered-eigenvalue bounds. Log-log fits use
the fixed windows 30--200, 60--400, and 120--800. These finite-resolution fits do
not prove an asymptotic exponent. The analytic optimal Sharpe is positive and
finite and is compared with a 100,000-period stock-return simulation (E6).

\subsection{Estimation, selection, and evaluation}
We use the paper's response-one ridge estimator in a controlled Nystr\"om
subspace of the actual Mat\'ern RKHS, without rescaling the kernel diagonal.
The inducing points are independent of all training and OOS samples. Empirical
effective complexity is
\[
\widehat{\mathcal C}(\lambda)=\operatorname{Tr}
\bigl[\widehat\Sigma_T(\widehat\Sigma_T+\lambda I)^{-1}\bigr].
\]
Training sample sizes and all 96 positive penalties are fixed before simulation.
Every replication generates genuine stock-level returns. Sample-size paths are
nested within a replication; different replications are independent. Validation
uses the last quarter of the available history, following an earlier training
block. After selecting lambda by validation loss, the policy is refitted on all
available observations. An independent 720-period block gives OOS loss and
unannualized Sharpe. Neither population moments, $W^\star$, nor OOS data enter
selection. Separate independent population quadrature integrates conditional
factor and idiosyncratic moments for more precise population-risk evaluation.

The population-risk oracle minimizes mean regret over the same fixed grid.
Boundary choices and all high-complexity observations are retained. We report
means, medians, pointwise Monte Carlo intervals, selection frequencies, and
benchmarks for strong shrinkage, near-unregularized policies, a validated affine
characteristic policy, and equal weights.
Selected-penalty rate regressions use Monte Carlo means. The selection panel in
the main complexity figure displays the median penalty; reporting both retains
the effect of skewness and boundary choices rather than conflating the two.

For each penalty, the exact decomposition in the numerical subspace is
\[
Q(\widehat W_\lambda)-Q(W^\star)
=\{Q(W_\lambda)-Q(W^\star)\}
+\|\widehat W_\lambda-W_\lambda\|_{\Sigma}^2
+2\langle\widehat W_\lambda-W_\lambda,
\Sigma W_\lambda-\mu\rangle.
\]
The cross term is signed; only the squared-norm estimation term is necessarily
nonnegative. Population bias includes the explicitly measured numerical
approximation floor. Rank sensitivity compares 128, 256, and 512 inducing points.
For numerical rank $r$, empirical complexity cannot exceed $\min(T,r)$.
The near-unregularized end of each curve is therefore a subspace diagnostic
and can be rank-sensitive. The paired acceptance check concerns oracle and
validation-selected regret; it is not a uniform approximation guarantee over
the entire penalty grid or an infinite-rank asymptotic experiment.

\subsection{Monte Carlo evidence and limits}
The headline uses 500 independent replications. At $T=1440$, mean oracle and validation-selected population losses are 0.988906 and 0.991938; their mean population Sharpes are 0.10756 and 0.09595. Mean empirical complexities are 8.227 and 10.341. The selected-to-oracle mean regret ratio is 1.972. The validation boundary-selection frequency is 8.4 percent. The ensemble oracle selects a grid boundary at 0 of 10 sample sizes. These observations are retained in all plots and rate regressions. Along the full largest-T grid, mean independent OOS loss is 0.999977 at the lowest represented complexity and 1.198229 at the highest. Its observed grid minimum is 0.988305 at mean empirical complexity 8.227. The observed mean OOS Sharpe maximum is 0.11088 at complexity 8.227. These OOS extrema are descriptive and never determine a fitted policy or a validation choice. The oracle-regret log-log slope is -0.518 over the full grid (whole-replication bootstrap interval [-0.535, -0.498]) and -0.692 over its predeclared upper half ([-0.742, -0.647]). The difference between these windows shows the finite-sample transition; neither window nor the calibration was selected to match the theoretical slope.

Rate regressions use the full T grid, its upper half, and its largest four values.
Bootstrap intervals resample entire independent replications, preserving their
joint T-by-lambda surfaces and recomputing both the oracle and validation choices.
The displayed theoretical slopes are benchmarks for the theory's risk envelope;
they are not imposed DGP-specific equalities for this fixed smooth population
policy. Slower convergence, boundary selections, and differences from benchmark
slopes are reported without modifying the calibration or fitting windows.

Robustness varies N independently of T, signal strength and idiosyncratic noise, characteristic persistence, and Mat\'ern smoothness. All 15 additional configurations use 200 replications. Nearby multiples of three, $N\in\{99,300,501,999\}$, preserve exact triplet balance. The weak/high-noise case uses $\mu_F/2$ and $\sigma_\varepsilon=0.12$; the strong/low-noise case uses $2\mu_F$ and $\sigma_\varepsilon=0.04$. Persistence is $\rho\in\{0.70,0.90,0.97\}$ and the alternative kernel smoothness values are $\nu\in\{0.5,2.5\}$ at the same dimension $D=6$. Empirical ranks are a separate finite-support experiment that violates the infinite-spectrum E5 condition. They use coordinatewise midranks $2(\operatorname{rank}(Z_i)-1/2)/N-1$ and retain E1--E4 and E6. Bounded heteroskedastic Gaussian shocks retain tail control but generally destroy the baseline conditional pricing identity; their conditional standard deviation is $\sigma_\varepsilon(z)=0.08\exp(0.35z_2)$. E1 and E2 remain exact, and bounded positive volatility preserves the E5 operator comparison with adjusted constants. E3 and E6 for the unrestricted heteroskedastic optimum are not asserted. Their regret is explicitly measured against the numerical population optimum in the same Nystr\"om subspace, not the baseline analytic optimum. The empirical JKP exercise is distinct from all of these artificial economies.

\input{\simoutput/tables/table1_calibration.tex}
\input{\simoutput/tables/table2_assumptions.tex}
\input{\simoutput/tables/table3_main_results.tex}
\input{\simoutput/tables/table4_rates.tex}
\input{\simoutput/tables/table5_robustness.tex}

\begin{figure}[p]
\centering
\includegraphics[width=\textwidth]{\simoutput/figures/baseline_four_panel.pdf}
\caption{Independent OOS performance against empirical effective complexity. All penalties are retained; loss uses a logarithmic scale.}
\label{fig:sim_baseline_four_panel}
\end{figure}

\begin{figure}[p]
\centering
\includegraphics[width=\textwidth]{\simoutput/figures/bias_estimation_regret.pdf}
\caption{Exact bias, signed estimation contribution, and total regret. A symmetric logarithmic scale retains negative cross terms.}
\label{fig:sim_bias_estimation_regret}
\end{figure}

\begin{figure}[p]
\centering
\includegraphics[width=\textwidth]{\simoutput/figures/population_spectrum.pdf}
\caption{Population kernel and managed spectra on common independent quadrature. Shaded fixed windows show spectral-fit sensitivity.}
\label{fig:sim_population_spectrum}
\end{figure}

\begin{figure}[p]
\centering
\includegraphics[width=\textwidth]{\simoutput/figures/rate_laws.pdf}
\caption{Oracle and validation rate paths. Dashed curves display theoretical slope benchmarks without constraining fitted results.}
\label{fig:sim_rate_laws}
\end{figure}

\begin{figure}[p]
\centering
\includegraphics[width=\textwidth]{\simoutput/figures/main_robustness.pdf}
\caption{Signal-to-noise robustness with the same fixed learning design.}
\label{fig:sim_main_robustness}
\end{figure}
